\chapter{Memory와 Metadata의 Area 최적화}

\section{사례 1: BRAM inference가 먼저다}

초기 synthesis는 true-dual-port memory의 두 write port를 합성기가 인식하는
형태로 표현하지 못해 RAM을 register로 풀었다. 네 bank hierarchy가 각각 약
27--29k LUT와 47,150 FF를 사용했고, 전체는 device 용량을 훨씬 초과했다.
이때의 우선순위는 address 폭 축소나 mux 미세 최적화가 아니라 inference
template 수정이다.

최종 \rtlfile{MDL\_XXX\_bram\_dp.v}는 port별 synchronous process를 분리한다.

\begin{lstlisting}
(* ram_style = "block" *) reg [DATA_W-1:0] mem [0:DEPTH-1];

always @(posedge iSYS_CLK) begin : PORT_A
  if (iENA) begin
    if (iWEA) mem[iADDRA] <= iDINA;
    oDOUTA <= mem[iADDRA];
  end
end

always @(posedge iSYS_CLK) begin : PORT_B
  if (iENB) begin
    if (iWEB) mem[iADDRB] <= iDINB;
    oDOUTB <= mem[iADDRB];
  end
end
\end{lstlisting}

\rtlfile{ram\_style} attribute만 붙이는 것으로 부족하다. port semantics가
primitive가 지원하는 동작과 맞아야 한다. read-first 동작과 same-row dual write가
없다는 scheduler 계약도 simulation assertion으로 남긴다.

\section{사례 2: logical depth와 physical depth를 분리한다}

최초 coefficient bank parameter는 11-bit address, depth 2048이었다. 최대 transform
길이는 1152이지만 네 physical bank가 각각 logical address 전체를 받을 수 있게
작성되어 있었다. 2048$\times$23 true-dual-port bank 하나는 7-series에서
RAMB36과 RAMB18의 조합으로 구현되었고, 네 bank가 6 BRAM-tile equivalents를
소비했다.

최종 mapping은 logical address $x$를 side와 row로 나눈다.
\[
row=x[10:1],\qquad side=x_0\oplus x_k.
\]
지원 transform에서 최대 row는 575이므로 각 bank는 576$\times$23이면 충분하다.
이 크기는 RAMB36E1 하나에 들어가며 네 coefficient bank는 총 4 RAMB36이 된다.

\subsection{왜 정보가 사라지지 않는가}

$row$에는 $x_k$가 남아 있다($k\ge1$). 그러므로
\[
x_0=side\oplus x_k
\]
로 제거한 bit 0을 복원할 수 있다. 즉 $(side,row)$는 logical address와 일대일
대응한다. 단순히 address LSB를 버린 것이 아니라, 그 정보를 side에 보존한
compact mapping이다.

\begin{lstlisting}
assign oSIDE       = addr[0] ^ addr[k];
assign oLOCAL_ADDR = addr[ADDR_W-1:1];
\end{lstlisting}

실제 RTL은 variable $k$ 대신 4, 6, 7의 고정 tap을 case로 선택한다. mapping별
bijection, row bound, radix-2 bank balance를 exhaustive test로 확인한다.

\section{사례 3: scheme별 ROM을 word packing으로 통합한다}

최초 \rtlfile{zeta\_rom}은 ML-KEM, ML-DSA, HAETAE, NTRU+768,
NTRU+864/1152 ROM을 모두 instantiate한 뒤 mode mux로 출력을 선택했다. 합성
결과의 twiddle storage는 RAMB18 세 개, 즉 1.5 tile이었다.

최종판은 2048$\times$32의 packed image 하나를 두 synchronous read port로
읽는다. 저장 구간은 다음과 같다.

\begin{longtable}{L{37mm}r r L{78mm}}
\toprule
Mode & Base word & Words & Packing \\
\midrule
ML-DSA & 0 & 256 & 23-bit value 1개/word \\
HAETAE & 256 & 128 & 16-bit value 2개/word \\
ML-KEM & 384 & 64 & 12-bit value 2개/word \\
NTRU+768 & 448 & 96 & 12-bit value 2개/word \\
NTRU+864/1152 & 544 & 144 & 12-bit value 2개/word \\
Falcon forward & 688 & 512 & 16-bit value 2개/word \\
Falcon inverse & 1200 & 512 & 16-bit value 2개/word \\
\bottomrule
\end{longtable}

총 1712 word이므로 2048-word memory에 들어간다. narrow mode의 logical address
LSB가 word 안의 low/high slot을 고른다.

\begin{lstlisting}
phys_w = base_w + logical_addr[10:1];
slot_w = logical_addr[0];

out = slot_q ? word_q[31:16] : word_q[15:0];
\end{lstlisting}

packed ROM은 두 RAMB36을 사용한다. twiddle memory만 보면 1.5 tile에서 2 tile로
0.5 tile 증가했지만, Falcon forward/inverse table까지 포함한다. coefficient
bank가 6 tile에서 4 tile로 줄어 전체 accelerator memory는 7.5 tile에서 6
tile로 1.5 tile 감소했다. 부분별 증감을 나눠 보아야 올바른 설명이 된다.

\section{사례 4: pipeline metadata도 area다}

data path만 세면 control storage를 놓치기 쉽다. 최초 core는 네 lane의
\rtlfile{va/vb/vc/vd}를 각각 12-stage pipeline으로 운반했고, BRAM read 뒤에
data/address/zeta를 다시 capture하는 register group도 있었다.

최종 transaction에서는 A/B/C lane이 항상 유효하고 radix-3일 때만 D가 없다.
따라서 tail valid는 layer-static butterfly mode로 계산한다.

\begin{lstlisting}
wire ctl_is_r3 = (ctl_but_sel_q == BUT_R3_NTT) ||
                 (ctl_but_sel_q == BUT_R3_INTT);
wire wb_lane_d_valid = !ctl_is_r3;
\end{lstlisting}

최종 13-stage 기준으로 네 lane-valid shift register 52 FF를 제거한다. 이 숫자는
해당 register의 구조적 감소량이며, 전체 FF delta와 동일시하지 않는다. 다른
pipeline과 interface register도 동시에 바뀌었기 때문이다.

read data와 synchronous zeta output은 butterfly C1 register로 직접 들어가도록
정렬하여 duplicate capture stage도 제거했다. pipeline register를 무조건 줄인
것이 아니라, 같은 cycle boundary가 두 번 존재하는 부분을 찾은 것이다.

\section{사례 5: 계측 logic을 production에서 생성하지 않는다}

LOAD/CORE/DUMP/TOTAL cycle counter는 개발 중 유용하지만 accelerator protocol의
필수 기능은 아니다. 최종 top은 parameterized generate로 네 32-bit counter와
incrementer를 production configuration에서 제거한다.

\begin{lstlisting}
generate
  if (ENABLE_PERF_COUNTERS != 0) begin : GEN_PERF_COUNTERS
    // four 32-bit cycle counters
  end else begin : GEN_NO_PERF_COUNTERS
    always @(posedge iSYS_CLK) begin
      oLOAD_CYCLES  <= 0;
      oCORE_CYCLES  <= 0;
      oDUMP_CYCLES  <= 0;
      oTOTAL_CYCLES <= 0;
    end
  end
endgenerate
\end{lstlisting}

하위 AXI load/dump beat counter도 같은 parameter를 사용한다. 검증용 관측성을
유지하면서 release hardware에 incrementer와 fanout을 남기지 않는 방법이다.

\section{Area 감소를 설명하는 올바른 문장}

최초 standalone과 최종 standalone의 역사적 결과는 LUT 5,278에서 4,736,
FF 3,072에서 2,792, BRAM tile 7.5에서 6.0, DSP 4에서 4다. 각각 10.3\%,
9.1\%, 20.0\% 감소와 DSP 동일이다. 그러나 package, top boundary, clock가
완전히 일치하는 단일 변수 실험은 아니다.

따라서 다음과 같이 쓴다.

\begin{goodbox}[권장 해석]
최종 implementation point는 최초 standalone archive보다 542 LUT, 280 FF,
1.5 BRAM tile을 적게 사용하고 같은 네 DSP를 사용한다. compact coefficient
mapping은 bank당 RAMB18을 제거했고, 다른 LUT/FF 변화에는 arithmetic,
scheduler, metadata, wrapper 차이가 함께 포함된다.
\end{goodbox}

특정 RTL 변경 하나가 전체 542 LUT를 줄였다고 쓰려면 그 변경만 다른 matched
run이 추가로 필요하다.
